% ===================== 第 11 章 =====================
\section{产品与工程实现}

\subsection{技术栈与事实来源}

后端 Python 3.12、FastAPI、Pydantic v2、SQLAlchemy 2.0 异步、Alembic、PostgreSQL 16（含 pgvector）；前端 React + TypeScript + React Flow，实时更新走 Server-Sent Events；测试 pytest。数据库是唯一事实来源：任务状态、议会进度、证据、审计记录全部落库，进程内存与工作流框架内部不保存任何唯一状态，Worker 重启不丢任务。

\subsection{模型网关：角色代码里没有厂商 SDK}

\code{OpenAICompatibleModelGateway} 对接任意 OpenAI Chat Completions 兼容端点。能力包括：流式解析并在失败时回退非流式；DeepSeek 思维链（\code{reasoning\_content}）单独成流、标注为过程材料；结构化输出修复至多一次；300 秒总时限；特定厂商参数冲突自动降级；token、费用、延迟、重试逐笔审计。议会侧只面向 \code{ModelGateway} 协议发送 \code{ModelRequest}，不感知厂商。用户输入的端点会先规范化（补全 scheme、剥离尾部斜杠、把已知控制台门户改写为 API 主机），保存前做一次真实连通探测；密钥不回显、不进错误信息。

\subsection{工具网关与取证管线}

\code{HttpToolGateway} 对接四个公开学术源：OpenAlex、Crossref、Semantic Scholar 免凭证可用，Unpaywall 按其要求携带联系邮箱。取证管线为一条可审计流水线：

\begin{lstlisting}
研究问题 -> 证据需求规划 -> 多源候选发现 -> 元数据归一化去重
-> 全文获取与结构化解析 -> 研究级/结果级发现抽取 -> 原文定位绑定
-> 设计质量与边界审查 -> 独立性审计 -> 引用蕴含验证
-> 事件账本 -> 证据门 -> 证据图
\end{lstlisting}

规范化作品（Canonical Work）按 DOI $\rightarrow$ PubMed/arXiv ID $\rightarrow$ 标题+一作+年份 $\rightarrow$ 模糊匹配的优先级，把预印本、正式版与扩展版识别为同一作品族，避免三个版本被计为三项独立研究。每一步留事件记录，因此任意结论都能回答「谁、在哪篇文献的哪一页、基于什么原文、以什么证据等级提出，审计员是否同意」。

\subsection{报告组装：两种格式，一次构建}

\code{ReportService} 从证据图、账本与任务行一次构建 Research Brief，Markdown（阅读与导出）与 JSON（界面渲染）来自同一次 build，结构上不可能对议会结论产生分歧；报告阶段不调用模型。最终论文只使用账本已记录的材料，模型失败时由确定性兜底组装，绝不输出空壳；任务终态不随论文生成成败改变，失败只记 \code{FINAL\_PAPER\_FAILED}。所有导出经 \code{sanitize\_export} 脱敏，签名 URL 与服务器本地路径不出现在交付物里。

\subsection{账户、知识库、技能与协作：产品层四模块}

\begin{table}[ht]
\centering\small
\caption{首版产品层模块（均为按账户隔离的持久能力）}
\label{tab:product}
\begin{tabularx}{\linewidth}{@{}L{2.8cm} X@{}}
\toprule
\rowcolor{accent}
\tblhead 模块 & \tblhead 能力与关键约束 \\
\midrule
账户体系 \code{accounts} & 邮箱验证码两阶段注册/找回；PBKDF2-HMAC-SHA256（20 万次迭代）口令；bearer 令牌只存 sha256；全部事实落库，API 进程重启不丢登录态；验证码发送/消耗均为单行原子操作 \\
个人知识库 \code{knowledge} & PDF/TXT/MD/CSV/DOCX/PPTX/XLSX 多格式入库时解析为文本，跨任务复用；Postgres 全文检索加 ILIKE 子串分支（后者保证中文可搜）；删除做引用检查，已成为证据的文档不可孤儿化 \\
研究者技能 \code{skills} & 从 GitHub 仓库解析安装 SKILL.md（常规路径优先、递归扫描兜底、API 限流量时下载整包扫描）；可选 LLM 辅助只做选择与摘要、不伪造文件；启用后按任务把技能文本注入议会提示词 \\
人工标注 \code{annotations} & 把待评盲点/主张冻结成标注批，多名标注者打标并计算 $\kappa$/$\alpha$ 一致性；标注是人类对系统产物的判断，永不进入证据图、投影器永不读取 \\
公开分享 \code{shared} & 终态任务可铸造 256 位随机只读分享链接；仅 GET、按令牌查询、快照经公开脱敏（费用、签名 URL、本地路径全部移除） \\
\bottomrule
\end{tabularx}
\end{table}

\subsection{前端工作台：中心是证据地图，不是聊天框}

工作台由 23 个视图组成，覆盖研究简报、七人议会实时状态、盲点雷达、争议证据地图（React Flow）、证据谱系、审计轨迹、演进视图、人工标注、知识库、技能管理、会话历史与账户设置。设计语言定位为科研仪器：颜色只表达证据状态，被反驳节点淡化但不隐藏，缺口计数常驻并逐条点名，不使用卡通形象、霓虹色与脉冲动画，完整支持 \code{prefers-reduced-motion}；结论与局限在同一视图并排出现。实时层做了有界化与批量合并（第 9.5 节），历史任务与后台切回场景不再因事件重放冻结界面。

\subsection{部署}

\code{docker-compose.yml} 编排六个服务：postgres、一次性 migrate（跑 Alembic 并建数据库角色）、api、worker、web（nginx 静态托管加反代）、caddy（唯一对外入口，域名就绪时一行配置自动签发 HTTPS）。GitHub Actions 在 push 到 main 后经 SSH 拉取、重建并做落地页健康检查；服务器 \code{.env} 与模型密钥由部署者维护，不进入仓库。

\subsection{安全与伦理}

系统输出不构成临床诊断或医疗建议，命中心理健康主题时自动附加科研辅助定位声明；不抓取未授权个人数据，不绕过付费墙；上传材料的存储、日志与导出均做隔离与脱敏；每份报告的局限一节固定声明 AI 辅助性质、证据覆盖范围与「不以模型置信度替代统计不确定性或专家判断」。
